%% =====================================================================
\section{The Transfer Principle}\label{sec:transfer}
%% =====================================================================

This section contains the load-bearing new statement of the paper. It is elementary; its
content is entirely in what it says about the relation between two constructions that the
literature has treated as incompatible.

\subsection{The coefficient functional}\label{ss:coeff}

\begin{definition}\label{def:rho}
Let $R\in\QQ(t)$ have poles only at non-positive integers and $\deg R\le-2$, and write its
partial-fraction decomposition as
\begin{equation}\label{eq:pf}
 R(t)=\sum_{i\ge1}\sum_{k\ge0}\frac{r_{i,k}}{(t+k)^i},\qquad r_{i,k}\in\QQ,
\end{equation}
the sums being finite. For $\theta$ in the domain of definition below, set
\begin{equation}\label{eq:rhodef}
 \rho_i:=\sum_k r_{i,k}\quad(i\ge1),
 \qquad
 \rho_{0,\theta}:=-\sum_{i,k}r_{i,k}\sum_{\nu=0}^{k-1}(\nu+\theta)^{-i},
\end{equation}
with the convention that the empty sum $\sum_{\nu=0}^{-1}$ is $0$. We call
\begin{equation}\label{eq:functional}
 \mathcal R:\ R\ \longmapsto\ (\rho_{0,\theta},\rho_2,\rho_3,\rho_4,\dots)
\end{equation}
the \emph{coefficient functional} of $R$ at $\theta$. It is manifestly $\QQ$-linear in the
partial-fraction data $(r_{i,k})$, and for fixed $\theta\in\QQ$ it takes values in $\QQ$.
\end{definition}

We note at once the vanishing that both worlds use.

\begin{lemma}[purity by degree]\label{lem:rho1}
\PROVED{} If $\deg R\le-2$ then $\rho_1=\lim_{t\to\infty}tR(t)=0$.
\end{lemma}

\begin{proof}
Multiplying \eqref{eq:pf} by $t$ and letting $t\to\infty$, every term with $i\ge2$
contributes $0$ and the terms with $i=1$ contribute $\sum_kr_{i,k}=\rho_1$; on the other
hand $tR(t)\to0$ because $\deg R\le-2$.
\end{proof}

\subsection{The principle}\label{ss:principle}

\begin{theorem}[Transfer Principle]\label{thm:transfer}
\PROVED{} Let $R$ be as in Definition~\ref{def:rho} and let $\widetilde R$ be a primitive of
$R$ (any antiderivative; the two statements below are insensitive to the constant, which is
killed by $\rho_1=0$).
\begin{enumerate}[label=\textup{(\roman*)},leftmargin=2.6em]
\item \textup{(archimedean)} For $\theta>0$,
\begin{equation}\label{eq:arch}
 \sum_{m\ge0}R(m+\theta)\;=\;\rho_{0,\theta}+\sum_{i\ge2}\rho_i\,\zeta(i,\theta).
\end{equation}
\item \textup{($p$-adic)} For $\theta\in\Qp$ with $|\theta|_p\ge q_p$ --- the domain on
which $\zeta_p(s,\theta)$ is defined \cite[\S2]{LaiSprang2023} ---
\begin{equation}\label{eq:padic}
 -\int_{\Zp}\widetilde R(t+\theta)\,\dd t
   \;=\;\rho_{0,\theta}+\sum_{i\ge2}\rho_i\,\omega(\theta)^{1-i}\zeta_p(i,\theta).
\end{equation}
\end{enumerate}
The rational numbers $\rho_{0,\theta}$ and $\rho_i$ occurring in \eqref{eq:arch} and in
\eqref{eq:padic} are \emph{the same}: both are the values of the single functional
\eqref{eq:functional}.
\end{theorem}

%% REFEREE [R2]: the hypothesis |theta|_p >= q_p is faithful to Lai-Sprang \S2, but
%% q_2 = 4 and |1/2|_2 = 2, so it EXCLUDES theta = 1/2 at p = 2 -- the integration shift
%% of LSZ, of Beukers' R^(B), and of the whole p = 2 cone of \S\ref{sec:scan}, at which
%% Proposition~\ref{prop:linform} nonetheless invokes "the case theta = theta_0 of
%% Theorem~\ref{thm:transfer}(ii)".  The substitute identity (LSZ Lemma 9) was already
%% quoted in the proof of Theorem~\ref{thm:classification}; it is now flagged here.
\begin{remark}[the excluded point $\theta=\tfrac12$, $p=2$]\label{rem:halfexcl}
Since $q_2=4$ and $|\tfrac12|_2=2$, the hypothesis of (ii) excludes $\theta=\tfrac12$ at
$p=2$ --- which is exactly the shift used by every $p=2$ construction below. There
$\zeta_2(s,\tfrac12)$ is not defined, and its role is played instead by
\cite[Lemma 9]{LSZ2025},
\[
 \frac1{s-1}\int_{\ZZ_2}\frac{\dd t}{(t+\tfrac12)^{s-1}}=2^s\,\zeta_2(s)
 \qquad(s\in\ZZ_{\ge2}),
\]
in which the single shift $\tfrac12$ already produces the full $\zeta_2$. Statement (i),
the coefficient functional \eqref{eq:functional} and both corollaries of
\S\ref{ss:corollaries} are unaffected, because they are statements about $\QQ$; what
changes at $(p,r)=(2,1)$ is only which $p$-adic \emph{value} is attached to $\rho_i$.
Everywhere else in the paper $|\theta_0|_p\ge q_p$ holds.
\end{remark}

\begin{proof}
Statement (ii) is \cite[Lemma 21]{LaiSprang2023}, whose proof we recall in outline because
it is what makes the comparison visible. Writing
$\widetilde R(t+\theta)=\sum_kr_{1,k}\log_p\langle t+k+\theta\rangle
 +\sum_{i\ge2}\sum_k r_{i,k}(1-i)^{-1}(t+k+\theta)^{1-i}$, one integrates termwise; the
translation formula for the Volkenborn integral \cite[Lemma 5]{LSZ2025} moves each shift
$k$ back to $0$ at the cost of the finite sum $\sum_{\nu<k}(\nu+\theta)^{-i}$, and the
identity $\zeta_p(i,\theta)=\omega(\theta)^{i-1}\cdot\frac{1}{i-1}\int_{\Zp}
 \dd t/(t+\theta)^{i-1}$ of \cite[\S2]{LaiSprang2023} converts the remaining integrals. The
$\log_p$ terms carry the factor $\rho_1$, which vanishes by Lemma~\ref{lem:rho1}. Collecting
the finite sums gives exactly $\rho_{0,\theta}$ of \eqref{eq:rhodef}.

Statement (i) is the classical Hurwitz expansion, obtained by the identical bookkeeping.
Summing \eqref{eq:pf} over $m\ge0$ termwise,
$\sum_{m\ge0}(m+k+\theta)^{-i}=\zeta(i,\theta)-\sum_{\nu=0}^{k-1}(\nu+\theta)^{-i}$ for
$i\ge2$, which produces $\sum_i\rho_i\zeta(i,\theta)$ together with precisely the finite
correction $\rho_{0,\theta}$; the $i=1$ terms again assemble into $\rho_1$ times a divergent
constant and vanish by Lemma~\ref{lem:rho1}, absolute convergence of the whole rearrangement
being guaranteed by $\deg R\le-2$.

The final sentence is not an additional claim but an observation about the two computations
just performed: in each, the coefficient of $\zeta(i,\cdot)$ resp.\ $\zeta_p(i,\cdot)$ is
$\sum_kr_{i,k}$, and the rational term is $-\sum_{i,k}r_{i,k}\sum_{\nu<k}(\nu+\theta)^{-i}$.
Neither expression refers to the completion.
\end{proof}

\begin{remark}\label{rem:sameness}
The two constructions differ \emph{only} in which value is attached to each $\rho_i$:
$\zeta(i,\theta)$ on one side, $\omega(\theta)^{1-i}\zeta_p(i,\theta)$ on the other. This is
the exact sense in which they are ``the same construction''. It is also, we note, the
precise content of the correspondence that LSZ point out in
\cite[Remark 13]{LSZ2025}, where the coefficients $\rho_{n,0},\rho_{n,3}$ of their $p$-adic
form are observed to be ``exactly the same'' as those of the archimedean expansion
$\sum_{m\ge0}R_n''(m+\tfrac12)=\rho_{n,0}+\rho_{n,3}(1-2^{-5})\zeta(5)$. What is new here is
not the observation at one point but its use as a \emph{principle}.
\end{remark}

\subsection{The two corollaries that do the work}\label{ss:corollaries}

\begin{corollary}[transfer of identities]\label{cor:transfer}
\PROVED{} Let $\mathcal F$ be any family of rational functions as in
Definition~\ref{def:rho}, parametrised by a set $P$, and let $g:P\to P$ be a map for which
one knows an identity of the form
\[
 \mathcal R(R_{g\mathbf a})=\kappa(\mathbf a)\cdot\mathcal R(R_{\mathbf a}),
 \qquad \kappa(\mathbf a)\in\QQ^{\times},
\]
i.e.\ a coefficientwise proportionality of the two coefficient vectors. Then that identity
holds verbatim for the archimedean linear forms \eqref{eq:arch} and for the Volkenborn
linear forms \eqref{eq:padic} at every admissible $\theta$, with the same $\kappa$.
\end{corollary}

\begin{proof}
Immediate from Theorem~\ref{thm:transfer}: both forms are the images of the same vector
$\mathcal R(R)$ under a fixed (completion-dependent) evaluation map, and the hypothesis is a
statement about $\mathcal R(R)$ alone.
\end{proof}

\begin{corollary}[transfer of denominators]\label{cor:den}
\PROVED{} Any statement of the form ``$\Lambda\cdot\rho_i\in\ZZ$ for all $i$'' or
``$\Lambda\cdot\rho_{0,\theta}\in\ZZ$'', with $\Lambda\in\QQ$ depending on the parameters
only, is a statement about rational numbers and therefore holds simultaneously for the
archimedean and the $p$-adic construction built from the same $R$.
\end{corollary}

Corollaries~\ref{cor:transfer} and \ref{cor:den} are the answer to the second LSZ question
of \S\ref{ss:questions}. The expected shape of an answer was a summation or transformation
theory for Volkenborn integrals. The actual answer is that none is required: the group acts
on $R$, the transformation identity is an identity between the rational vectors
$\mathcal R(R)$, and the Volkenborn integral is a passive spectator. In the terminology of
the campaign this is the $p$-adic analogue of the $\nu_p$-sharpening principle, pushed one
level further --- not merely ``valid in every completion'' but ``the $p$-adic
construction's coefficients \emph{are} the archimedean construction's coefficients''.

\begin{remark}\label{rem:whatitdoesnotgive}
Three things Theorem~\ref{thm:transfer} does \emph{not} give, stated here so that the reader
is not misled by \S\ref{sec:group}. (a) It does not say that the $p$-adic \emph{smallness}
transfers; smallness is a statement about the value, not the coefficients, and it is exactly
where the two worlds diverge (\S\ref{sec:ledger}). (b) It does not say that admissibility
transfers; a rational function whose numerator bricks sit at the wrong shift is a perfectly
good archimedean object and a bad $p$-adic one (Proposition~\ref{prop:F2}). (c) It does not
give a group where there is none: the family LSZ actually use has trivial transfer structure
(Finding~\ref{find:F3}).
\end{remark}

\subsection{The savings function is completion-independent}\label{ss:deltainsens}

For the Rhin--Viola machinery one needs to know that the denominator saving is the same
object in both worlds. It is, and for a cheap reason.

%% REFEREE [R5]: the proof as drafted did not establish o(n).  Its own ingredients give
%% O(n/log n) primes times O(log ell) = O(log n) each = O(n) -- precisely what has to be
%% excluded -- and the displayed chain "O(n) O(log n)/log n o(1)" is not a well-formed
%% bound.  (The same defect is in DIG_GROUP.md \S1b, which writes the change as
%% O(pi(m_3 n)), i.e. O(1) per prime rather than O(log ell).)  Evidence class lowered
%% from [PROVED] to [DERIVED]; the argument replaced by the direction-perturbation one,
%% with the crude bound stated honestly.  Nothing downstream depends on this: the
%% proposition can only help a conclusion that is negative anyway.
\begin{proposition}\label{prop:deltainsens}
\DERIVED{} Let the parameters of a hypergeometric brick be integers $c\cdot n$ with $c$ a
fixed direction vector, and let $\Phi_n=\prod_{\sqrt{m_0n}<\ell\le m_3n}\ell^{\nu_\ell}$ be
the Rhin--Viola/Zudilin denominator saving, $\nu_\ell=\max_{g\in\Gp}\ord_\ell(\Pi(cn)/\Pi(gcn))$.
Then, replacing each integer parameter $c\cdot n$ by a $\theta$-shifted parameter
$c\cdot n+O(1)$ leaves
\[
 \dgr:=\lim_{n\to\infty}\frac{\log\Phi_n}{n}
\]
unchanged, so that $\dgr$ is the same function of the direction in the $p$-adic and in the
archimedean setting.
\end{proposition}

\begin{proof}[Argument]
An $O(1)$ shift changes $\ord_\ell$ of each factorial by at most $O(1)$ and hence each
$\nu_\ell$ by $O(1)$. Bounding each of the $\pi(m_3n)-\pi(\sqrt{m_0n})=O(n/\log n)$ primes
in the window by its own $O(\log\ell)=O(\log n)$ gives only $O(n)$, which is not enough; the
point is that the shift moves $\lfloor c_in/\ell\rfloor$ only when a carry occurs, i.e.\ for
a set of primes of relative density $O(1/\ell)$, so that the expected total displacement is
$O\bigl(\sum_{\ell\le m_3n}\log\ell/\ell\bigr)=O(\log n)=o(n)$. Equivalently and more
structurally: with the shift absorbed into the direction, the perturbed direction is
$c+O(1/n)\to c$, and $\dgr$ is given by the integrals \eqref{eq:deltadef} of the piecewise
constant $\Lambda$, whose discontinuity set is finite and $\dd\psi$-null away from $0$; the
integrals are therefore continuous along $c+O(1/n)\to c$. We do not carry the second
argument out in detail, and label the statement \DERIVED{} accordingly.
\end{proof}

Proposition~\ref{prop:deltainsens} is what licenses computing $\dgr$ entirely with the
classical archimedean machinery, which is what \S\ref{sec:group} does.
